At A the exterior accessible sector has angle $240^\circ$ and radius three, with area $6\pi$. Wrapping around either neighbouring vertex B or F leaves length two and adds a $60^\circ$ sector, totalling $\dfrac{4\pi}{3}$. Reaching C via B leaves length one and adds a $60^\circ$ sector of area $\dfrac{\pi}{6}$. Reaching E via F also leaves one, but the missing edge DE allows a $180^\circ$ sector rather than the usual $60^\circ$, adding $\dfrac{\pi}{2}$. The sectors have disjoint interiors; reaching D along the other boundary uses the entire leash. Thus total area is $6\pi+\dfrac{4\pi}{3}+\dfrac{\pi}{6}+\dfrac{\pi}{2}=\boxed{8\pi}$.
\begin{asy}
import graph;import geometry;usepackage("amsmath");size(7cm);
pair[] V;for(int i=0;i<6;++i)V.push(dir(60*i));for(int i=0;i<6;++i)if(i!=3)draw(V[i]--V[(i+1)%6],linewidth(1.5));draw(V[3]--V[4],dashed+red);string[] nm={"A","B","C","D","E","F"};for(int i=0;i<6;++i)label("$"+nm[i]+"$",V[i],unit(V[i]));draw(V[0]--V[5]--V[4],blue);label("removed edge",(V[3]+V[4])/2,(-1,-1));
\end{asy}

